%% title: Light-trap images, the out-of-domain dataset
%% target: Sections/Methods.tex
%% anchor: \subsection{Integration of arbitrary encoders (backbones)}\label{sec:projection}
%% action: before
%% why: Adds the third dataset as its own subsection under Datasets. Fixes which export of the trap images the paper uses: the draft's tables use the older export, where 13.7 % of crops carry keys outside the vocabulary, mostly from key mismatch.
%% source: data/flemming (47,905 crops, 486 species, 12 traps, nights 2022-06-16 to 2023-11-07); journal/2026-07-28-flemming-generalization.md for the two exports
%% check: The AMI hardware reference (REF). | Should we say who identified the crops (Flemming Helsing) here or in the acknowledgements? | Asger's AMI numbers must be re-run on this export before they can sit next to ours.
%% status: pasted
\subsubsection{Light-trap images}\label{sec:data_traps}
The out-of-domain data are crops of insects photographed by twelve Automated Moth Identification (AMI) light traps \textbf{(REF)}, four at each of three sites in Denmark, on nights between June 2022 and November 2023.
Each crop was identified to species by an expert.
The set holds 47,905 crops of 486 species.
Species names were matched to GBIF keys by name, after which only 234 crops of six species fall outside the training vocabulary.
The crops differ from the training images in camera, illumination, background and pose, and each moth is photographed many times per night by the traps' time-lapse.
